% ============================================================================
%  CrowdROM Experiments: Design, Math, and Results
% ============================================================================
\documentclass[11pt]{article}

% ---- Packages ----
\usepackage[margin=1in]{geometry}
\usepackage{amsmath, amssymb, amsthm}
\usepackage{booktabs}
\usepackage{multirow}
\usepackage{hyperref}
\usepackage{xcolor}
\usepackage{enumitem}
\usepackage[font=small]{caption}
\usepackage{subcaption}
\usepackage{float}
\usepackage{array}
\usepackage{colortbl}
\usepackage{listings}
\usepackage{tikz}
\usepackage{pgfplots}
\pgfplotsset{compat=1.18}
\usepackage{siunitx}

% ---- Code listings ----
\definecolor{codebg}{gray}{0.96}
\definecolor{codegreen}{rgb}{0.1,0.5,0.1}
\definecolor{codeblue}{rgb}{0.1,0.1,0.6}
\definecolor{codegray}{rgb}{0.4,0.4,0.4}
\lstset{
  language=Python,
  basicstyle=\ttfamily\small,
  backgroundcolor=\color{codebg},
  keywordstyle=\color{codeblue}\bfseries,
  commentstyle=\color{codegreen}\itshape,
  numbers=left,
  numberstyle=\tiny\color{codegray},
  numbersep=5pt,
  frame=single,
  breaklines=true,
  showstringspaces=false,
  tabsize=4,
  xleftmargin=2em,
  framexleftmargin=1.5em,
  captionpos=b,
}

% ---- Colors ----
\definecolor{goodgreen}{HTML}{2E7D32}
\definecolor{badred}{HTML}{C62828}
\definecolor{warnbrown}{HTML}{E65100}
\definecolor{lightgray}{HTML}{F5F5F5}
\definecolor{accentblue}{HTML}{1565C0}

% ---- Custom commands ----
\newcommand{\R}{\mathbb{R}}
\newcommand{\bx}{\mathbf{x}}
\newcommand{\ba}{\mathbf{a}}
\newcommand{\by}{\mathbf{y}}
\newcommand{\bU}{\mathbf{U}}
\newcommand{\bA}{\mathbf{A}}
\newcommand{\bp}{\mathbf{p}}
\newcommand{\bmu}{\boldsymbol{\mu}}
\newcommand{\brho}{\boldsymbol{\rho}}
\newcommand{\Rtwo}{R^2}
\newcommand{\good}[1]{\textcolor{goodgreen}{\textbf{#1}}}
\newcommand{\bad}[1]{\textcolor{badred}{#1}}
\newcommand{\warn}[1]{\textcolor{warnbrown}{#1}}
\newcommand{\regime}[1]{\textsf{#1}}
\DeclareMathOperator*{\argmin}{arg\,min}
\DeclareMathOperator*{\argmax}{arg\,max}

\title{%
  \textbf{CrowdROM Experimental Log}\\[6pt]
  \large Complete Experiment Design, Mathematical Formulation,\\
  \large and Results for All Suites (DYN through VFIX)}
\author{Macias-Olivares, E.}
\date{February--March 2026}

\begin{document}
\maketitle
\tableofcontents
\newpage


% ============================================================================
\section{Problem Setting and Physics}
\label{sec:physics}
% ============================================================================

\subsection{Vicsek--Morse Particle Model}

We study $N = 200$ interacting agents on a doubly-periodic domain
$\Omega = [0, L_x) \times [0, L_y)$ with $L_x = L_y = 25$.
Each particle $i$ at discrete time $k$ has position
$\bx_i^k \in \Omega$ and heading angle $\theta_i^k$.

\paragraph{Update rules.}
The heading is updated via Vicsek alignment plus Morse forces:
\begin{equation}\label{eq:heading}
  \theta_i^{k+1}
  = \arg\!\biggl(\,
    \sum_{j:\,|\bx_j^k - \bx_i^k| < R}
    e^{i\theta_j^k}\biggr)
  + \eta\,\xi_i^k
  + \alpha_{\text{morse}}\,\Theta_i^{\text{Morse}, k},
\end{equation}
where $\eta$ is noise intensity, $\xi_i^k \sim \text{Uniform}(-\pi, \pi)$,
and $\Theta_i^{\text{Morse},k}$ is the heading correction from
the Morse potential gradient.

The position is updated at constant speed $v_0$:
\begin{equation}
  \bx_i^{k+1} = \bx_i^k + v_0\,\Delta t
  \begin{pmatrix}\cos\theta_i^{k+1}\\\sin\theta_i^{k+1}\end{pmatrix}
  \;\mod\;(L_x, L_y).
\end{equation}

\paragraph{Morse interaction forces.}
The pairwise Morse potential between particles $i$ and $j$:
\begin{equation}
  W(r) = -C_r\,e^{-r/l_r} + C_a\,e^{-r/l_a},
  \qquad r = |\bx_i - \bx_j|,
\end{equation}
produces short-range repulsion (strength $C_r$, range $l_r$) and
long-range attraction (strength $C_a$, range $l_a$).
The net force on particle $i$ is:
\begin{equation}
  \mathbf{F}_i^{\text{Morse}} = -\sum_{j \ne i} \nabla_{\bx_i} W(|\bx_i - \bx_j|).
\end{equation}


\subsection{The Seven Dynamical Regimes}

All experiments span 7 physically distinct regimes that cover the
parameter space:

\begin{table}[H]
\centering
\caption{The 7 dynamical regimes used throughout all experiments.}
\label{tab:regimes}
\small
\begin{tabular}{@{}clccccl@{}}
\toprule
\# & Regime & Speed $v_0$ & $\eta$ & $C_a$ & $C_r$ & Character \\
\midrule
1 & \regime{Gentle}       & 1.5  & 0.1 & 0.8  & 0.3  & Slow, smooth, easy \\
2 & \regime{HyperVelocity}& 15.0 & 0.2 & 1.5  & 0.5  & Extreme speeds \\
3 & \regime{HyperNoisy}   & 4.0  & 1.5 & 1.5  & 0.5  & Near-random noise \\
4 & \regime{Blackhole}    & 4.0  & 0.2 & 12.0 & 0.1  & Extreme attraction \\
5 & \regime{Supernova}    & 4.0  & 0.2 & 0.1  & 10.0 & Extreme repulsion \\
6 & \regime{Baseline}     & 4.0  & 0.2 & 1.5  & 0.5  & Balanced reference \\
7 & \regime{PureVicsek}   & 4.0  & 0.2 & ---  & ---  & No Morse forces \\
\bottomrule
\end{tabular}
\end{table}


% ============================================================================
\section{ROM Pipeline Mathematics}
\label{sec:rom_math}
% ============================================================================

\subsection{Step 1: Density Field via KDE}

Particle positions are converted to a continuum field:
\begin{equation}
  \rho^k(x,y) = (G_\sigma * H^k)(x,y),
  \qquad
  H^k_{jl} = \frac{1}{\Delta x\,\Delta y}
  \sum_{i=1}^N \mathbf{1}[\bx_i^k \in C_{jl}],
\end{equation}
where $H^k$ is a $64 \times 64$ histogram and $G_\sigma$ is a Gaussian
kernel with bandwidth $\sigma = 5$ grid cells, applied with periodic
wrapping.  Each frame $\brho^k \in \R^{4096}$ ($= 64 \times 64$).


\subsection{Step 2: Shift Alignment (sPOD)}

Before POD, we factor out translational motion via FFT cross-correlation:
\begin{equation}
  (\delta_y^k, \delta_x^k)
  = \argmax_{(s_y, s_x)}
  \text{Re}\bigl[\mathcal{F}^{-1}(\hat{\rho}_{\text{ref}} \odot \overline{\hat{\rho}^k})\bigr]_{s_y, s_x},
\end{equation}
where $\rho_{\text{ref}} = \frac{1}{T}\sum_k \rho^k$ is the temporal mean.
The aligned frame is:
\begin{equation}
  \tilde{\rho}^k(j,l) = \rho^k\bigl((j-\delta_y^k) \bmod N_y,\;
  (l - \delta_x^k) \bmod N_x\bigr).
\end{equation}


\subsection{Step 3: Density Transform (Optional)}

Two transform pipelines were tested:

\paragraph{Suite A ($\sqrt{}$+simplex+C0).}
\begin{enumerate}[nosep]
  \item Square-root transform: $\tilde{\rho} \leftarrow \sqrt{\tilde{\rho}}$
    (Hellinger embedding, stabilises variance),
  \item C0 clamping: clip negatives to zero after each rollout step,
  \item Simplex mass projection: $\hat{\rho} \leftarrow \hat{\rho} \cdot M_0 / \sum\hat{\rho}$.
\end{enumerate}

\paragraph{Suite B (raw+C2).}
\begin{enumerate}[nosep]
  \item Identity (no transform),
  \item C2 clamping: soft ReLU with small $\epsilon > 0$,
  \item No mass projection.
\end{enumerate}


\subsection{Step 4: POD Compression}

We compute the snapshot SVD:
\begin{equation}
  \tilde{\mathbf{X}} = \mathbf{U}\,\boldsymbol{\Sigma}\,\mathbf{V}^\top,
  \qquad
  \tilde{\mathbf{X}} \in \R^{MT \times n},
\end{equation}
and retain $d = 19$ modes capturing $> 99\%$ of variance:
\begin{equation}
  \ba^k = \mathbf{U}_d^\top\,\tilde{\brho}^k \in \R^d.
\end{equation}
The POD coefficient trajectory $\{\ba^k\}_{k=1}^T$ is what the ROM
methods forecast.


\subsection{Step 5a: MVAR Forecasting}
\label{sec:mvar}

The Multivariate Autoregressive (MVAR) model of order $p$:
\begin{equation}\label{eq:mvar}
  \hat{\ba}_{t+1} = \sum_{j=1}^p \mathbf{A}_j\,\ba_{t-j+1} + \mathbf{c},
\end{equation}
where $\mathbf{A}_j \in \R^{d \times d}$ and $\mathbf{c} \in \R^d$.

Writing $\mathbf{Z}_t = [\ba_t^\top, \ba_{t-1}^\top, \ldots, \ba_{t-p+1}^\top, 1]^\top \in \R^{dp+1}$
and $\boldsymbol{\Theta} = [\mathbf{A}_1, \ldots, \mathbf{A}_p, \mathbf{c}] \in \R^{d \times (dp+1)}$,
the model is $\hat{\ba}_{t+1} = \boldsymbol{\Theta}\,\mathbf{Z}_t$.

\paragraph{Ridge regression fit:}
\begin{equation}
  \hat{\boldsymbol{\Theta}} = \argmin_{\boldsymbol{\Theta}}
  \sum_t \|\ba_{t+1} - \boldsymbol{\Theta}\mathbf{Z}_t\|^2
  + \alpha\,\|\boldsymbol{\Theta}\|_F^2,
\end{equation}
with $\alpha = 10^{-4}$.  In closed form:
\begin{equation}
  \hat{\boldsymbol{\Theta}} = \mathbf{Y}\mathbf{Z}^\top(\mathbf{Z}\mathbf{Z}^\top + \alpha\mathbf{I})^{-1}.
\end{equation}


\subsection{Step 5b: LSTM Forecasting}
\label{sec:lstm}

\paragraph{Architecture.}
The \texttt{LatentLSTMROM} processes sequences of POD coefficients:
\begin{align}
  \mathbf{h}_t, \mathbf{c}_t &= \text{LSTM}(\mathbf{y}_t, \mathbf{h}_{t-1}, \mathbf{c}_{t-1}), \\
  \tilde{\mathbf{h}} &= \text{LayerNorm}(\mathbf{h}_L), \\
  \hat{\mathbf{y}}_{L+1} &= W_o\,\tilde{\mathbf{h}} + \mathbf{b}_o,
\end{align}
where $\mathbf{y}_t \in \R^d$ ($d=19$), $\mathbf{h}_t \in \R^{N_h}$,
and $L$ is the context window length.

\paragraph{LSTM gates.}
The standard LSTM cell computes:
\begin{align}
  \mathbf{f}_t &= \sigma(\mathbf{W}_{if}\mathbf{y}_t + \mathbf{W}_{hf}\mathbf{h}_{t-1} + \mathbf{b}_f), \\
  \mathbf{i}_t &= \sigma(\mathbf{W}_{ii}\mathbf{y}_t + \mathbf{W}_{hi}\mathbf{h}_{t-1} + \mathbf{b}_i), \\
  \tilde{\mathbf{c}}_t &= \tanh(\mathbf{W}_{ig}\mathbf{y}_t + \mathbf{W}_{hg}\mathbf{h}_{t-1} + \mathbf{b}_g), \\
  \mathbf{c}_t &= \mathbf{f}_t \odot \mathbf{c}_{t-1} + \mathbf{i}_t \odot \tilde{\mathbf{c}}_t, \\
  \mathbf{o}_t &= \sigma(\mathbf{W}_{io}\mathbf{y}_t + \mathbf{W}_{ho}\mathbf{h}_{t-1} + \mathbf{b}_o), \\
  \mathbf{h}_t &= \mathbf{o}_t \odot \tanh(\mathbf{c}_t).
\end{align}

\paragraph{Residual connection (optional).}
When enabled:
$\hat{\mathbf{y}}_{t+1} = \mathbf{y}_t + \Delta\hat{\mathbf{y}}$,
where $\Delta\hat{\mathbf{y}} = W_o\,\tilde{\mathbf{h}} + \mathbf{b}_o$.

\paragraph{Parameter count.}
For $L_{\text{layers}}$ LSTM layers, LayerNorm, and linear output:
\begin{equation}
  N_{\text{params}} = \underbrace{4(d + N_h + 1)N_h}_{\text{first layer}}
  + \underbrace{4(N_h + N_h + 1)N_h \cdot (L_{\text{layers}}-1)}_{\text{additional layers}}
  + \underbrace{2N_h}_{\text{LayerNorm}}
  + \underbrace{(N_h + 1)d}_{\text{output}}.
\end{equation}

\begin{table}[H]
\centering
\caption{Parameter counts for $d = 19$.}
\label{tab:params}
\small
\begin{tabular}{@{}ccr@{}}
\toprule
$N_h$ & $L_{\text{layers}}$ & Params \\
\midrule
16  & 1 & 2,723 \\
32  & 1 & 7,475 \\
64  & 1 & 23,123 \\
128 & 1 & 78,995 \\
128 & 2 & 211,475 \\
32  & 2 & 15,923 \\
64  & 2 & 56,403 \\
\bottomrule
\end{tabular}
\end{table}


\paragraph{Training procedure.}
\begin{itemize}[nosep]
  \item \textbf{Input normalization:} per-mode z-scoring using training-set $\mu, \sigma$.
  \item \textbf{Optimizer:} Adam, weight decay $10^{-5}$, learning rate $10^{-3}$.
  \item \textbf{LR schedule:} ReduceLROnPlateau (factor 0.5, patience 20).
  \item \textbf{Gradient clipping:} max norm $= 5.0$.
  \item \textbf{Early stopping:} patience 40 epochs (monitor val loss).
  \item \textbf{Max epochs:} 200.
  \item \textbf{Batch size:} 32.
  \item \textbf{Context window:} $p = 5$ time steps.
\end{itemize}


\subsection{Step 5c: WSINDy PDE Discovery (On Full Fields)}
\label{sec:wsindy_exp}

WSINDy discovers the PDE governing the 3-field system
$(\rho, p_x, p_y)$ directly, then time-integrates the discovered
equations.  The full mathematical detail is in the companion document
\texttt{wsindy\_math\_and\_code.tex}; here we summarise the key
equations:

\paragraph{Discovered system.}
\begin{align}
  \rho_t &= c_1\,\nabla\cdot\bp + c_2\,\Delta\rho
    + c_3\,\nabla\cdot(\rho\nabla\Phi) + c_4\,\Delta(\rho^2) + \ldots \\
  (p_x)_t &= d_1\,p_x + d_2\,|\bp|^2 p_x + d_3\,\partial_x\rho
    + d_4\,\Delta p_x + d_5\,\rho\,\partial_x\Phi + \ldots \\
  (p_y)_t &= e_1\,p_y + e_2\,|\bp|^2 p_y + e_3\,\partial_y\rho
    + e_4\,\Delta p_y + e_5\,\rho\,\partial_y\Phi + \ldots
\end{align}
Coefficients $c_m, d_m, e_m$ are identified by weak-form regression
(MSTLS, see companion document).

\paragraph{Weak $\Rtwo$.}
The quality of the discovered PDE is measured by:
\begin{equation}
  \Rtwo_{\text{weak}} = 1 - \frac{\|\mathbf{G}_{\text{active}}\mathbf{w} - \mathbf{b}\|^2}{\|\mathbf{b}\|^2}.
\end{equation}

\paragraph{Time integration.}
Discovered PDEs are integrated using ETDRK4 (if stiff linear terms
present) or RK4, with mass conservation and non-negativity projections.


\subsection{Step 6: Reconstruction and Evaluation}

The predicted density is reconstructed via:
\begin{equation}
  \hat{\rho}^k_{\text{phys}} = \text{un-align}\bigl(
    \text{inverse-transform}\bigl(
      \mathbf{U}_d\,\hat{\ba}^k + \bar{\brho}
    \bigr),\; \hat{\boldsymbol{\delta}}^k
  \bigr).
\end{equation}

\paragraph{Metrics.}
\begin{itemize}[nosep]
  \item \textbf{Rollout $\Rtwo$:} Coefficient of determination across all forecast
    pixels and times.  Autoregressive (no teacher forcing):
    \begin{equation}
      \Rtwo_{\text{rollout}} = 1 - \frac{\sum_{k,j,l}(\rho^k_{jl} - \hat\rho^k_{jl})^2}
      {\sum_{k,j,l}(\rho^k_{jl} - \bar\rho_{jl})^2}.
    \end{equation}
  \item \textbf{1-step $\Rtwo$:} Teacher-forced (ground-truth input at each step).
  \item \textbf{Negativity fraction:} $\%$ of cells with $\hat\rho < 0$.
  \item \textbf{POD ceiling:} $\Rtwo$ if the ``prediction'' is just reconstructing
    ground truth via the POD basis (upper bound for any ROM method).
\end{itemize}


% ============================================================================
\section{Experiment Suite VDYN: MVAR vs.\ LSTM Baseline}
\label{sec:vdyn}
% ============================================================================

\subsection{Design}
\begin{itemize}[nosep]
  \item \textbf{7 jobs:} one per dynamical regime.
  \item \textbf{Methods:} MVAR ($p=5$, $\alpha=10^{-4}$) and LSTM ($N_h=128$, $L=2$,
    211k params, residual=ON, $\sqrt{}$+simplex+C0).
  \item \textbf{POD:} $d=19$ modes, shift-aligned.
  \item \textbf{Forecast horizon:} 50\,s autoregressive rollout.
\end{itemize}

\subsection{Results}

\begin{table}[H]
\centering
\caption{VDYN: MVAR vs.\ LSTM rollout $\Rtwo$ (50\,s forecast).}
\label{tab:vdyn}
\small
\begin{tabular}{@{}l S[table-format=+1.3] S[table-format=+1.3] S[table-format=2.1] @{}}
\toprule
Regime & {MVAR $\Rtwo$} & {LSTM $\Rtwo$} & {LSTM Neg\%} \\
\midrule
\regime{Gentle}       & +0.904 & -1.90 & 18.3 \\
\regime{HyperVelocity}& +0.850 & -5.64 & 26.1 \\
\regime{HyperNoisy}   & +0.783 & -0.68 & 12.3 \\
\regime{Blackhole}    & +0.836 & -2.73 & 22.7 \\
\regime{Supernova}    & +0.734 & -3.22 & 19.4 \\
\regime{Baseline}     & +0.880 & -4.91 & 23.8 \\
\regime{PureVicsek}   & +0.643 & +0.12 & 11.4 \\
\midrule
\textbf{Mean}         & \textbf{+0.804} & \textbf{$-$2.71} & \textbf{19.1} \\
\bottomrule
\end{tabular}
\end{table}

\textbf{Key finding:} LSTM achieves strong 1-step $\Rtwo \approx 0.95$ yet
produces catastrophically negative rollout $\Rtwo$ in 6/7~regimes.
The problem is purely autoregressive error compounding.


% ============================================================================
\section{Experiment Suite VLST: Density Transform Ablation}
\label{sec:vlst}
% ============================================================================

\subsection{Design}

\textbf{Question:} Does the density transform pipeline affect LSTM rollout stability?

\begin{itemize}[nosep]
  \item \textbf{14 jobs:} 2 suites $\times$ 7 regimes.
  \item \textbf{Architecture:} Same as VDYN ($N_h=128$, $L=2$, 211k params, residual=ON).
  \item \textbf{Suite A:} $\sqrt{}$ + simplex mass projection + C0 clamping.
  \item \textbf{Suite B:} Raw (identity) transform + C2 clamping.
  \item Oscar job IDs via commit \texttt{ec1d8d0}.
\end{itemize}

\subsection{Results}

\begin{table}[H]
\centering
\caption{VLST: rollout $\Rtwo$ by transform pipeline.}
\label{tab:vlst}
\small
\begin{tabular}{@{}l S[table-format=+1.2] S[table-format=+1.2] S[table-format=2.1] S[table-format=2.1] @{}}
\toprule
Regime & {Suite~A ($\sqrt{}$+S)} & {Suite~B (raw)} & {Neg\% A} & {Neg\% B} \\
\midrule
\regime{Gentle}       & -1.37 & -0.41 & 5.2  & 14.2 \\
\regime{HyperVelocity}& -5.47 & -1.29 & 11.2 & 23.8 \\
\regime{HyperNoisy}   & -0.63 & -0.19 & 4.1  & 9.7  \\
\regime{Blackhole}    & -3.82 & -1.14 & 8.7  & 21.4 \\
\regime{Supernova}    & -2.91 & -0.88 & 6.3  & 17.3 \\
\regime{Baseline}     & -3.10 & -1.49 & 9.8  & 22.1 \\
\regime{PureVicsek}   & -1.81 & -0.72 & 8.4  & 22.3 \\
\midrule
\textbf{Mean}         & \textbf{$-$2.73} & \textbf{$-$0.87} & \textbf{7.7} & \textbf{18.7} \\
\bottomrule
\end{tabular}
\end{table}

\subsection{Findings}
\begin{itemize}[nosep]
  \item Raw+C2 wins in 6/7 regimes for rollout $\Rtwo$ (mean $-0.87$ vs.\ $-2.73$).
  \item However, raw+C2 produces \emph{more} negative densities (18.7\% vs.\ 7.7\%).
  \item \textbf{Neither transform achieves positive $\Rtwo$}: transform alone
    cannot fix 211k-param overfitting with residual connection.
  \item The $\sqrt{}$ transform amplifies small errors near zero, exacerbating divergence.
\end{itemize}


% ============================================================================
\section{Experiment Suite VARCH: Architecture Ablation}
\label{sec:varch}
% ============================================================================

\subsection{Design}

\textbf{Question:} Which architectural choice causes LSTM rollout divergence?

\begin{itemize}[nosep]
  \item \textbf{42 jobs:} 6 variants $\times$ 7 regimes.
  \item \textbf{All use:} $\sqrt{}$+simplex+C0 transform.
  \item Oscar job IDs via commit \texttt{92fc67d}.
\end{itemize}

\begin{table}[H]
\centering
\caption{VARCH variant definitions (all use $\sqrt{}$+simplex+C0, dropout=0).}
\label{tab:varch-design}
\small
\begin{tabular}{@{}cccccr@{}}
\toprule
ID & $N_h$ & $L$ & Residual & LayerNorm & Params \\
\midrule
A16   & 16 & 1 & ON  & ON  & 2,723  \\
A32   & 32 & 1 & ON  & ON  & 7,475  \\
A32x2 & 32 & 2 & ON  & ON  & 15,923 \\
\midrule
\textbf{B1}   & 32 & 1 & \good{OFF} & ON  & 7,475  \\
B2    & 32 & 1 & ON  & OFF & 7,475  \\
B3    & 32 & 1 & ON  & ON  & 7,475  \\
\bottomrule
\end{tabular}
\end{table}

B3 adds multi-step loss (4-step horizon during training).

\subsection{Capacity Sweep Results (A16, A32, A32x2)}

All capacity variants (with residual=ON) produce negative rollout $\Rtwo$:

\begin{center}
\begin{tabular}{lS[table-format=+1.2]l}
\toprule
Variant & {Mean $\Rtwo$} & Verdict \\
\midrule
A16 (2.7k params) & -0.63 & Negative \\
A32 (7.5k params) & -0.51 & Negative \\
A32x2 (15.9k params) & -0.68 & Negative \\
\bottomrule
\end{tabular}
\end{center}

\textbf{Conclusion:} Right-sizing the network (from 211k to 2.7k--16k params)
without removing the residual connection does \emph{not} fix rollout divergence.
Capacity is a \textbf{red herring}.


\subsection{Trick Ablation Results (B1, B2, B3)}

\begin{table}[H]
\centering
\caption{VARCH rollout $\Rtwo$ --- full 7-regime results.
\good{Green}: positive $\Rtwo$ (only B1 achieves this).}
\label{tab:varch-r2}
\small
\setlength{\tabcolsep}{4pt}
\begin{tabular}{@{}l *{7}{S[table-format=+1.2]} S[table-format=+1.2] @{}}
\toprule
Variant & {Gentle} & {HVel} & {HNoisy} & {BHole} & {SNova} & {Base} & {PVicsek} & {Mean} \\
\midrule
A16      & -0.62 & -0.73 & -0.28 & -0.91 & -0.55 & -0.89 & -0.41 & -0.63 \\
A32 (REF)& -0.48 & -0.58 & -0.33 & -0.72 & -0.44 & -0.71 & -0.30 & -0.51 \\
A32x2    & -0.81 & -0.69 & -0.38 & -0.95 & -0.61 & -0.84 & -0.45 & -0.68 \\
\rowcolor{lightgray}
\textbf{B1 (no res.)} & -0.03 & \good{+0.05} & \good{+0.24} & \good{+0.33} & -0.08 & -0.11 & -0.04 & \good{+0.05} \\
B2 (no LN)  & -0.79 & -1.03 & -0.59 & -1.12 & -0.85 & -1.07 & -0.62 & -0.87 \\
B3 (+multi) & -0.67 & -0.81 & -0.42 & -0.99 & -0.58 & -0.93 & -0.44 & -0.69 \\
\midrule
MVAR     & +0.90 & +0.85 & +0.78 & +0.84 & +0.73 & +0.88 & +0.64 & +0.80 \\
\bottomrule
\end{tabular}
\end{table}


% ============================================================================
\section{Root Cause Analysis}
\label{sec:rootcause}
% ============================================================================

\subsection{Finding \#1: Residual Connection is the Primary Killer}

The residual (delta) formulation:
\begin{equation}
  \hat{\mathbf{y}}_{t+1} = \mathbf{y}_t + \Delta\hat{\mathbf{y}}_t,
  \quad
  \Delta\hat{\mathbf{y}}_t = f_\theta(\mathbf{y}_{t-p+1:t}).
\end{equation}

During autoregressive rollout, prediction errors $\epsilon_t$ in
$\Delta\hat{\mathbf{y}}$ accumulate \emph{additively}:
\begin{equation}
  \hat{\mathbf{y}}_T = \mathbf{y}_0 + \sum_{t=0}^{T-1}(\Delta\mathbf{y}_t^* + \epsilon_t)
  = \mathbf{y}_T^* + \sum_{t=0}^{T-1}\epsilon_t.
\end{equation}
The cumulative error grows as $O(\sqrt{T})$ (random walk) or worse.
With $T \sim 500$ rollout steps, even $\epsilon \sim 10^{-3}$ yields
$O(10^{-2}\text{--}10^{-1})$ drift.

\textbf{Evidence:} Removing the residual (B1 vs.\ A32):
\begin{itemize}[nosep]
  \item Mean $\Rtwo$: $-0.51 \to +0.05$ (\good{$+$53 points}).
  \item B1 wins 7/7 regimes.
  \item B1 is the \emph{only} variant achieving positive rollout $\Rtwo$
    (3 out of 7 regimes).
  \item Only 4/56 experiments achieved positive rollout $\Rtwo$;
    all 4 are from~B1.
\end{itemize}


\subsection{Finding \#2: LayerNorm Helps, Multi-Step Hurts}

Relative to A32 reference ($\overline{\Rtwo} = -0.51$):
\begin{center}
\begin{tabular}{lS[table-format=+1.2]l}
\toprule
Ablation & {$\Delta\overline{\Rtwo}$} & Interpretation \\
\midrule
B1 (no residual) & +0.56 & Dominant improvement \\
A16 (tiny)        & -0.12 & Capacity: slight degradation \\
A32x2 (deeper)    & -0.17 & Depth: slight degradation \\
B3 (+multistep)   & -0.18 & Multistep hurts \\
B2 (no LayerNorm) & -0.36 & LayerNorm important \\
\bottomrule
\end{tabular}
\end{center}

\begin{figure}[H]
\centering
\begin{tikzpicture}
\begin{axis}[
    xbar,
    width=0.75\textwidth,
    height=5.5cm,
    xlabel={$\Delta\overline{R^2}$ vs.\ A32 reference},
    symbolic y coords={
      {B3: +multistep},
      {B2: no LayerNorm},
      {A32x2: deeper},
      {A16: tiny},
      {B1: no residual}
    },
    ytick=data,
    xmin=-0.5,
    xmax=0.65,
    bar width=10pt,
    nodes near coords,
    nodes near coords align={horizontal},
    every node near coord/.append style={font=\footnotesize},
    grid=major,
    major grid style={dotted,gray!50},
]
\addplot[fill=accentblue!60, draw=accentblue] coordinates {
    (-0.18,{B3: +multistep})
    (-0.36,{B2: no LayerNorm})
    (-0.17,{A32x2: deeper})
    (-0.12,{A16: tiny})
    (+0.56,{B1: no residual})
};
\draw[dashed, red, thick] (0,{rel axis cs:0,0}) -- (0,{rel axis cs:0,1});
\end{axis}
\end{tikzpicture}
\caption{Effect size of each ablation vs.\ A32 reference ($\overline{\Rtwo}=-0.51$).}
\label{fig:effect-sizes}
\end{figure}


\subsection{Finding \#3: 1-Step vs.\ Rollout Disconnect}

All variants achieve 1-step $\Rtwo \approx 0.95$ (teacher-forced).
The bottleneck is exclusively autoregressive error compounding.
This confirms that \emph{exposure bias} is the failure mechanism,
not per-step prediction quality.


% ============================================================================
\section{Experiment Suite VFIX: Combining All Winning Traits}
\label{sec:vfix}
% ============================================================================

\subsection{Design Rationale}

Based on root-cause analysis, VFIX combines:
\begin{itemize}[nosep]
  \item \textbf{Residual = OFF} (the biggest single improvement, $+53$ pts).
  \item \textbf{LayerNorm = ON} (stabilises training, $+36$ pts vs.\ OFF).
  \item \textbf{Raw + C2 transform} (better $\Rtwo$ than $\sqrt{}$+simplex in VLST).
  \item \textbf{Multistep = OFF} (hurt in prior experiments).
  \item \textbf{Dropout = 0.0} (no regularisation benefit observed).
\end{itemize}

\subsection{VFIX Variant Design}

\textbf{56 jobs} = 8 variants $\times$ 7 regimes.  Oscar job IDs 562954--563014,
commit \texttt{b430514}.

\begin{table}[H]
\centering
\caption{VFIX variants. All use residual=OFF, LayerNorm=ON, dropout=0, multistep=OFF.}
\label{tab:vfix-design}
\small
\begin{tabular}{@{}cllccclr@{}}
\toprule
Suite & ID & Purpose & $N_h$ & $L$ & Transform & Special & Params \\
\midrule
\multirow{4}{*}{A}
  & F16  & Alvarez-scale    & 16  & 1 & raw+C2 & ---  & 2,723  \\
  & F32  & Sweet spot (REF) & 32  & 1 & raw+C2 & ---  & 7,475  \\
  & F64  & Wider            & 64  & 1 & raw+C2 & ---  & 23,123 \\
  & F128 & Wide             & 128 & 1 & raw+C2 & ---  & 78,995 \\
\midrule
\multirow{2}{*}{B}
  & F32x2 & Deeper          & 32  & 2 & raw+C2 & ---  & 15,923 \\
  & F64x2 & Deep + wide     & 64  & 2 & raw+C2 & ---  & 56,403 \\
\midrule
\multirow{2}{*}{C}
  & FSS   & Sched.\ sampling & 32  & 1 & raw+C2 & SS=ON & 7,475 \\
  & FSQRT & Transform ctrl   & 32  & 1 & $\sqrt{}$+S+C0 & --- & 7,475 \\
\bottomrule
\end{tabular}
\end{table}

\subsection{Key Questions Addressed}

\begin{enumerate}[nosep]
  \item \textbf{Width sweep (F16--F128):} With residual=OFF, does wider help or hurt?
    Can Alvarez-scale ($\sim$2.7k params) match larger models?
  \item \textbf{Depth sweep (F32x2, F64x2):} Does depth provide benefit when the
    residual connection is already removed?
  \item \textbf{Scheduled sampling (FSS):} Does curriculum-based self-feeding reduce
    exposure bias further when residual is already gone?
  \item \textbf{Transform control (FSQRT):} With all fixes applied, which transform
    pipeline wins: raw+C2 or $\sqrt{}$+simplex+C0?
\end{enumerate}


% ============================================================================
\section{Early DYN Experiments: MVAR/LSTM/WSINDy Comparison}
\label{sec:dyn}
% ============================================================================

\subsection{DYN1--DYN7: Initial 3-Method Suite}

The initial DYN suite ran MVAR, LSTM, and WSINDy PDE discovery across
all 7 regimes.  The WSINDy component operated on the full
$(\rho, p_x, p_y)$ field system, not in latent space.

\begin{itemize}[nosep]
  \item \textbf{MVAR:} $p=5$, ridge $\alpha=10^{-4}$, in POD latent space.
  \item \textbf{LSTM:} $N_h=128$, $L=2$, 211k params, residual=ON.
  \item \textbf{WSINDy:} Full multi-field discovery with Morse library.
\end{itemize}

MVAR succeeded in all 7 regimes.  LSTM diverged in 6/7.
The original DYN LSTM jobs failed due to a backward-compatibility
bug in config loading (old keys \texttt{hidden\_dim}, \texttt{n\_layers})
which was later fixed via \texttt{\_COMPAT\_MAP} in \texttt{lstm\_rom.py}.


\subsection{XABL1--XABL8: Preprocessing Factorial Design}

An earlier ablation suite tested $2^3$ factorial over three preprocessing
factors:
\begin{center}
\begin{tabular}{lccc}
\toprule
Factor & Level 0 & Level 1 \\
\midrule
Density transform & raw & $\sqrt{}$ \\
Mass correction & none & simplex projection \\
Spectral scaling & OFF & ON (re-weight POD modes) \\
\bottomrule
\end{tabular}
\end{center}

\textbf{Key finding:} Shift alignment dominates all preprocessing
choices.  $\sqrt{}$ + simplex help 1--3\% in $\Rtwo$ but become
negligible once alignment is applied.  Alignment alone raises $\Rtwo$
from $\sim 0.5\text{--}0.7$ to $\sim 0.95+$.


% ============================================================================
\section{Summary of All Experiments}
\label{sec:summary}
% ============================================================================

\subsection{Experiment Inventory}

\begin{table}[H]
\centering
\caption{Complete experiment inventory.}
\small
\begin{tabular}{@{}llrrl@{}}
\toprule
Suite & Purpose & Jobs & Commit & Status \\
\midrule
XABL & Preprocessing $2^3$ factorial & 8 & --- & Completed \\
DYN1--7 & 3-method comparison & 7 & \texttt{658c59b} & Completed (LSTM failed) \\
VLST & Transform ablation & 14 & \texttt{ec1d8d0} & Completed \\
VARCH & Architecture ablation & 42 & \texttt{92fc67d} & Completed \\
VFIX & Fix suite (all winning traits) & 56 & \texttt{b430514} & Submitted \\
\midrule
\textbf{Total} & & \textbf{127} & & \\
\bottomrule
\end{tabular}
\end{table}


\subsection{Key Conclusions Across All Suites}

\begin{enumerate}
  \item \textbf{MVAR is the strongest ROM method} (mean $\Rtwo = +0.804$)
    across all 7 regimes.  Simple, fast, reliable.

  \item \textbf{The LSTM residual connection is the dominant failure mode}
    for autoregressive rollout.  Removing it yields the only positive
    rollout $\Rtwo$ ($+53$-point improvement).

  \item \textbf{Model capacity is a red herring.}  Downsizing from 211k
    to 2.7k parameters does not fix the residual-driven accumulation.

  \item \textbf{LayerNorm is important:} removing it degrades $\Rtwo$ by
    $-36$ points.  Multi-step loss hurts ($-18$ points).

  \item \textbf{The density transform matters secondarily.}  Raw+C2
    outperforms $\sqrt{}$+simplex for rollout $\Rtwo$ but at the cost
    of more negative densities.

  \item \textbf{Shift alignment is the single most impactful preprocessing
    step}, providing $+40$--$50$ points of $\Rtwo$ improvement for all methods.

  \item \textbf{WSINDy operates on the full PDE system} and is
    complementary to the latent-space methods (MVAR/LSTM).  It provides
    interpretable governing equations rather than black-box forecasts.
\end{enumerate}


\subsection{The MVAR--LSTM Gap}

\begin{center}
\begin{tabular}{lS[table-format=+1.2]l}
\toprule
Method & {Mean $\Rtwo$} & Status \\
\midrule
MVAR (VDYN) & +0.80 & Dominant \\
LSTM best (VARCH-B1, no residual) & +0.05 & Gap: $-75$ pts \\
LSTM worst (VDYN, 211k+residual) & -2.71 & Catastrophic \\
\bottomrule
\end{tabular}
\end{center}

VFIX aims to close this gap below 50 points by combining the winning
architectural and transform choices.


% ============================================================================
\section{Compute Resources and Infrastructure}
\label{sec:compute}
% ============================================================================

\subsection{Oscar HPC Cluster (Brown University)}
\begin{itemize}[nosep]
  \item \textbf{Access:} \texttt{ssh emaciaso@ssh.ccv.brown.edu}
  \item \textbf{Per job:} 4 CPU cores, 32\,GB RAM, 12\,h wall time
  \item \textbf{Storage:} \texttt{/oscar/scratch/emaciaso/oscar\_output} (512\,GB quota)
  \item \textbf{Code sync:} Git-based (local $\to$ GitHub $\to$ Oscar \texttt{git pull})
\end{itemize}

\subsection{Software Engineering}
\begin{itemize}[nosep]
  \item \textbf{Config generator:} Python scripts generating YAML configs per experiment.
  \item \textbf{SLURM orchestrator:} Batch scripts submitting all jobs with logging.
  \item \textbf{Analysis pipeline:} \texttt{analyze\_all\_ablations.py} parses JSON summaries,
    builds comparison tables, computes effect sizes.
  \item \textbf{Backward-compat fix:} \texttt{\_COMPAT\_MAP} in \texttt{lstm\_rom.py}
    translates old DYN config keys to new format.
\end{itemize}

\subsection{Total Compute Budget}

\begin{center}
\begin{tabular}{lrrr}
\toprule
Suite & Jobs & Max hours/job & Total core-hours \\
\midrule
XABL & 8 & 12 & 384 \\
DYN & 7 & 12 & 336 \\
VLST & 14 & 12 & 672 \\
VARCH & 42 & 12 & 2,016 \\
VFIX & 56 & 12 & 2,688 \\
\midrule
\textbf{Total} & \textbf{127} & & \textbf{$\sim$6,096} \\
\bottomrule
\end{tabular}
\end{center}


\end{document}
